\section{Conclusion}\label{sec:Conclusion}

Teacher bias is typically assessed by examining  gaps in teacher's subjective evaluations, conditional on objective measures of student abilities. A key challenge in identifying such bias is the presence of measurement error in test scores, which are used as proxies for student ability. We show that both classical and certain forms of non-classical measurement error introduce equivalent bias in the estimated conditional gap in subjective evaluations. We discuss three different strategies to address this contamination bias. While an IV approach is more commonly used in applied econometric studies, we demonstrate that an EIV approach can address measurement error under weaker assumptions.

The empirical analysis focuses on conditional SES gaps in teacher track recommendations in the Netherlands. This is a setting where potential teacher bias is highly consequential and findings on conditional gaps with respect to parental SES are widely accepted. Using new administrative data, we find that measurement error in test scores can explain a substantial portion of the conditional SES gap, between 35 and 43\%.

It may seem striking that measurement error explains such a large share of the conditional SES gap, especially since our test scores are reliable by conventional norms. Yet when student ability is strongly correlated with SES and matters for teacher track recommendations, contamination bias can still be substantial. While our findings show that teacher bias may be smaller than previously reported, the strong correlation between ability and SES points to significant inequalities of opportunity before and during primary education.